\documentclass[10pt]{article}
% Shared LaTeX style for MIT 18.06 exam revision notes.
% Compile topic files with Tectonic, XeLaTeX, or LuaLaTeX.

\usepackage[a4paper,margin=0.72in]{geometry}
\usepackage{fontspec}
\usepackage{amsmath,amssymb,mathtools}
\usepackage{booktabs,array,tabularx}
\usepackage{enumitem}
\usepackage{titlesec}
\usepackage{fancyhdr}
\usepackage{microtype}
\usepackage{xcolor}

\setmainfont{Times New Roman}
\setsansfont{Arial}

\setlength{\parindent}{0pt}
\setlength{\parskip}{3.6pt}
\setlength{\abovedisplayskip}{6pt}
\setlength{\belowdisplayskip}{6pt}
\setlength{\abovedisplayshortskip}{4pt}
\setlength{\belowdisplayshortskip}{4pt}
\renewcommand{\arraystretch}{1.08}
\setlength{\tabcolsep}{4.5pt}

\titleformat{\section}
  {\normalfont\large\bfseries}
  {\thesection}
  {0.55em}
  {}
\titlespacing*{\section}{0pt}{10pt}{3pt}

\titleformat{\subsection}
  {\normalfont\normalsize\bfseries}
  {\thesubsection}
  {0.5em}
  {}
\titlespacing*{\subsection}{0pt}{7pt}{2pt}

\setlist[itemize]{leftmargin=1.35em,itemsep=1pt,topsep=2pt,parsep=0pt}
\setlist[enumerate]{leftmargin=1.55em,itemsep=1pt,topsep=2pt,parsep=0pt}

\pagestyle{fancy}
\fancyhf{}
\fancyfoot[C]{\scriptsize\color{black!45}MIT 18.06 Linear Algebra Exam Notes --- \thepage}
\renewcommand{\headrulewidth}{0pt}
\renewcommand{\footrulewidth}{0pt}

\newcommand{\notesubtitle}[1]{%
  {\centering\itshape #1\par}
  \vspace{2pt}
}

\newcommand{\source}[1]{%
  {\centering\small #1\par}
  \vspace{6pt}
}

\newcommand{\labelnote}[2]{%
  \par\smallskip
  \noindent\textbf{#1.}\quad #2
  \par\smallskip
}

\newcommand{\tightlabel}[1]{%
  \par\smallskip
  \noindent\textbf{#1.}\quad
}

\newcolumntype{Y}{>{\raggedright\arraybackslash}X}
\newcolumntype{L}[1]{>{\raggedright\arraybackslash}p{#1}}

\newenvironment{notetable}
  {\small\begin{tabularx}{\linewidth}}
  {\end{tabularx}}


\begin{document}

\begin{center}
{\LARGE 3.4--3.5 Independence, Basis, Dimension, and Subspaces\par}
\end{center}
\notesubtitle{Concise exam revision notes for independence, span, bases, and subspace dimensions.}
\source{Based on handwritten MIT 18.06 revision notes.}

\section{Core Idea}

\labelnote{Core idea}{Independence, span, and basis are all controlled by the equation \(Ax=0\), where the columns of \(A\) are the vectors being tested.}

For vectors \(v_1,\ldots,v_n\), put them into a matrix
\[
A=\begin{bmatrix}v_1&v_2&\cdots&v_n\end{bmatrix}.
\]
Then the dependence test is
\[
x_1v_1+x_2v_2+\cdots+x_nv_n=0
\qquad\Longleftrightarrow\qquad
Ax=0.
\]
The vectors are independent exactly when the only solution is \(x=0\).

\section{Independence and Dependence}

\begin{center}
\small
\begin{tabularx}{\linewidth}{@{}L{0.28\linewidth}Y Y@{}}
\toprule
Situation & Meaning & Exam consequence \\
\midrule
Independent columns & \(Ax=0\) has only \(x=0\). & No free variables. \\
Dependent columns & \(Ax=0\) has a nonzero solution. & At least one free variable. \\
More columns than rows & \(n>m\). & Columns must be dependent. \\
\bottomrule
\end{tabularx}
\end{center}

\labelnote{Exam trigger}{To test whether vectors are independent, solve \(Ax=0\), not \(Ax=b\).}

If \(Ax=0\) has a nonzero solution, that solution gives an explicit linear relation among the columns. For example,
\[
x_1v_1+x_2v_2+x_3v_3=0
\]
with some nonzero coefficient means one vector is forced by the others.

\section{Span}

The span is the set of all linear combinations of the given vectors:
\[
\operatorname{span}\{v_1,\ldots,v_k\}
=\{c_1v_1+\cdots+c_kv_k\}.
\]
It answers the question: which vectors can be produced?

\begin{center}
\small
\begin{tabularx}{0.82\linewidth}{@{}L{0.34\linewidth}Y@{}}
\toprule
Vectors in \(\mathbb{R}^2\) & Span \\
\midrule
\((1,0)^{T},\ (0,1)^{T}\) & All of \(\mathbb{R}^2\). \\
\((1,0)^{T},\ (-1,0)^{T}\) & Only the \(x\)-axis. \\
\bottomrule
\end{tabularx}
\end{center}

\labelnote{Common mistake}{Having two vectors in \(\mathbb{R}^2\) does not automatically mean they span \(\mathbb{R}^2\). They must point in independent directions.}

\section{Basis}

A basis is a set of vectors that is both independent and spanning.

\begin{center}
\small
\begin{tabularx}{\linewidth}{@{}L{0.25\linewidth}Y Y@{}}
\toprule
Property & What it prevents & What it guarantees \\
\midrule
Independent & No repeated direction. & No vector is unnecessary. \\
Spanning & No missing direction. & Every vector in the space is reachable. \\
Basis & Both at once. & Coordinates are unique. \\
\bottomrule
\end{tabularx}
\end{center}

\labelnote{Exam memory}{A basis is the smallest spanning set and the largest independent set.}

\section{Row Space and Column Space}

For an \(m\times n\) matrix \(A\), the column vectors have length \(m\), while row vectors have length \(n\).

\begin{center}
\small
\begin{tabularx}{\linewidth}{@{}L{0.23\linewidth}Y L{0.18\linewidth}Y@{}}
\toprule
Space & Built from & Lives in & Basis rule \\
\midrule
Column space \(C(A)\) & Columns of \(A\). & \(\mathbb{R}^m\) & Use pivot columns of the original \(A\). \\
Row space \(C(A^{T})\) & Rows of \(A\). & \(\mathbb{R}^n\) & Use nonzero rows of echelon form \(R\). \\
\bottomrule
\end{tabularx}
\end{center}

Example:
\[
A=
\begin{bmatrix}
2&4\\
3&6
\end{bmatrix},
\qquad
R=
\begin{bmatrix}
1&2\\
0&0
\end{bmatrix}.
\]
The row space can be represented from \(R\):
\[
C(A^{T})=\operatorname{span}\{(1,2)\}.
\]
The column space must use the pivot column from the original matrix:
\[
C(A)=\operatorname{span}\left\{
\begin{bmatrix}2\\3\end{bmatrix}
\right\}.
\]

\labelnote{Common mistake}{Do not use the pivot columns of \(R\) as a basis for \(C(A)\). Row operations change the columns, but they preserve the row space.}

\section{Dimension and Rank}

The dimension of a subspace is the number of vectors in a basis for that subspace.

\begin{center}
\small
\begin{tabularx}{\linewidth}{@{}L{0.24\linewidth}L{0.22\linewidth}Y@{}}
\toprule
Quantity & Value & Meaning \\
\midrule
\(\dim C(A)\) & \(r\) & Number of pivot columns. \\
\(\dim C(A^{T})\) & \(r\) & Number of nonzero rows in \(R\). \\
\(\dim N(A)\) & \(n-r\) & Number of free variables. \\
\(\dim N(A^{T})\) & \(m-r\) & Number of missing pivots among rows. \\
\bottomrule
\end{tabularx}
\end{center}

The row space and column space have the same dimension:
\[
\dim C(A)=\dim C(A^{T})=r.
\]
They usually live in different spaces, so they are not the same subspace.

\section{Four Fundamental Subspaces}

For \(A\in\mathbb{R}^{m\times n}\), the four subspaces are:

\begin{center}
\small
\begin{tabularx}{\linewidth}{@{}L{0.24\linewidth}L{0.18\linewidth}L{0.18\linewidth}Y@{}}
\toprule
Subspace & Lives in & Dimension & How to find a basis \\
\midrule
Column space \(C(A)\) & \(\mathbb{R}^m\) & \(r\) & Pivot columns of \(A\). \\
Row space \(C(A^{T})\) & \(\mathbb{R}^n\) & \(r\) & Nonzero rows of \(R\). \\
Nullspace \(N(A)\) & \(\mathbb{R}^n\) & \(n-r\) & Special solutions of \(Ax=0\). \\
Left nullspace \(N(A^{T})\) & \(\mathbb{R}^m\) & \(m-r\) & Special solutions of \(A^{T}y=0\). \\
\bottomrule
\end{tabularx}
\end{center}

\labelnote{Dimension check}{The dimensions split by ambient space: \(r+(n-r)=n\) in \(\mathbb{R}^n\), and \(r+(m-r)=m\) in \(\mathbb{R}^m\).}

\section{Rank-One Reminder}

A rank-one matrix is one column times one row:
\[
A=uv^{T}.
\]
For example,
\[
\begin{bmatrix}
a\\b\\c
\end{bmatrix}
\begin{bmatrix}
2&3&7&8
\end{bmatrix}
\]
has rank \(1\) if both factors are nonzero.

\labelnote{Exam use}{Rank-one matrices are useful because all columns are multiples of one vector and all rows are multiples of one row.}

\section{Common Mistakes}

\begin{center}
\small
\begin{tabularx}{\linewidth}{@{}L{0.40\linewidth}Y@{}}
\toprule
Mistake & Correct exam habit \\
\midrule
Testing independence with \(Ax=b\). & Test with \(Ax=0\). \\
Forgetting that \(n>m\) forces dependence. & More vectors than dimension means at least one free variable. \\
Using columns of \(R\) for \(C(A)\). & Use pivot columns of original \(A\). \\
Saying row space and column space are the same. & They have the same dimension \(r\), but may live in different spaces. \\
Counting all columns as a basis. & Count pivot columns only. \\
\bottomrule
\end{tabularx}
\end{center}

\section{Final Exam Checklist}

\tightlabel{Checklist}
\begin{enumerate}
  \item Write the size: \(A\) is \(m\times n\).
  \item Put the vectors as columns if testing independence or span.
  \item Solve \(Ax=0\) to test independence.
  \item Find the rank \(r\) from pivots.
  \item Use pivot columns of \(A\) for a basis of \(C(A)\).
  \item Use nonzero rows of \(R\) for a basis of \(C(A^{T})\).
  \item Use special solutions for a basis of \(N(A)\).
  \item Check dimensions: \(r\), \(r\), \(n-r\), and \(m-r\).
  \item State whether the vectors are independent, spanning, a basis, or dependent.
\end{enumerate}

\end{document}
